\devoirDS{1}{Chapitres 1 à 3}{Identifier un liquide/6,Une préparation précise/7,Diluer avant de doser/7}
\begin{blocLivre}{ds-1-1}\begin{exoBac}{6}{Identifier un liquide}
Le récipient vide pèse 35,2 g ; rempli de 25,0 mL d’un liquide, il pèse 54,9 g. Références à la même température : eau 1,00 ; éthanol 0,789 en g/mL.
\begin{questions}
\item Calculer la masse et la masse volumique du liquide.
\item Proposer une identification et en indiquer une limite.
\item Quatre nouvelles mesures donnent 0,79 ; 0,80 ; 0,78 ; 0,79 g/mL. Calculer la moyenne puis proposer une amélioration du protocole.
\end{questions}
\end{exoBac}\end{blocLivre}
\begin{blocLivre}{ds-1-2}\begin{exoBac}{7}{Une préparation précise}
On veut préparer 200,0 mL d’une solution de soluté à \qty{12,0}{g.L^{-1}} par dissolution.
\schema{dissolution}
\begin{questions}
\item Identifier le solvant si la solution est aqueuse.
\item Calculer la masse de soluté.
\item Expliquer les étapes et le choix de verrerie.
\item Le trait est dépassé : expliquer l’effet sur la concentration et pourquoi retirer une partie de la solution ne répare pas l’erreur.
\end{questions}
\end{exoBac}\end{blocLivre}
\begin{blocLivre}{ds-1-3}\begin{exoBac}{7}{Diluer avant de doser}
On prélève 20,0 mL de solution inconnue et on complète à 100,0 mL. La fille a une masse volumique de \qty{1,030}{g.mL^{-1}} selon la gamme fournie.
\schema{etalonnage}
\begin{questions}
\item Lire la concentration de la fille.
\item Calculer le facteur de dilution et la concentration mère.
\item Quelle verrerie permet le prélèvement ?
\item Pourquoi une lecture hors gamme demanderait-elle une nouvelle préparation ?
\end{questions}
\end{exoBac}\end{blocLivre}

